\textbf{V-axis encoding peaks in posterior visual cortex across all
five frequency bands.}
\textbf{(a--e)} MNE-style scalp topomaps, one per band ($\delta$,
$\theta$, $\alpha$, $\beta$, $\gamma$), of the cohort Pearson
correlation between the LLM V-axis ($28$ stimulus projections) and
per-channel band power on FACED ($n{=}123$ subjects). Diverging
colourmap: positive (red) vs.\ negative (blue) correlation; symmetric
limits across bands. Gold circles mark the top channel per band;
text under each topomap names that channel and its $r$ value
(strongest cell PO3/$\gamma$, $r{=}{+}0.48$). The colourbar (centred)
gives Pearson $r$; the small head diagram (right) shows the
$32$-channel 10--20 montage used in all five panels (A1/A2 mastoids
excluded from plotting). \textbf{(f)} Region-mean $|r|$ aggregated
over all five bands: occipital ($0.21$) $>$ parietal ($0.18$) $>$
central ($0.18$) $>$ frontal ($0.16$). The cohort signal is
posterior-dominant for video-evoked emotion---a paradigm-level
finding for dynamic stimuli rather than a refutation of the static
frontal-alpha asymmetry literature, which we replicate qualitatively
in direction at smaller magnitude.
